\section{Recognizing points on residual curves}\label{sec:residual-curves}

Corollary~\ref{cor:flag-matching} identifies the residue character spaces at
quasi-prime \((d-1)\)-divisors. These residue fields have transcendence
degree one, but their point valuations are not yet distinguished inside
their character spaces. We now recognize those terminal residues whose
constant fields are algebraic closures of finite fields. For these
residues we also recognize every point inertia line. The construction
uses Pop's density theorem for an individual field, followed by reduction
of its characters modulo \(\ell\).

Fix a function field \(F/\kappa\) over an algebraically closed field
of characteristic different from \(\ell\), of transcendence degree
\(d\geq 2\), and let
\(v\) be a quasi-prime \((d-1)\)-divisor. Put
\[
 F_0=Fv,\qquad \kappa_0=\kappa v.
\]
By Lemma~\ref{lem:residue-duality}, restriction to units induces a
surjective map \(\rho_v:D_v\to A_{F_0}\) with kernel \(I_v\).
Define
\begin{equation}\label{eq:residual-inertia-set}
 S_F^{\mathrm{in}}
   =\overline{\bigcup_{w\ {\rm quasi\text{-}prime\ divisor}} I_w}
       \subset A_F,
 \qquad
 J_v=\rho_v(S_F^{\mathrm{in}}\cap D_v)\subset A_{F_0}.
\end{equation}
Here the bar denotes the closure of the union, rather than the subgroup
generated by that union. In particular \(J_v\) is a set of characters,
not in general a subspace. It is stable under multiplication by
\(\Lambda\), so it makes sense to consider the family of lines contained
in \(J_v\). Both sets in \eqref{eq:residual-inertia-set}, and this family
of lines, are preserved by the isomorphisms already constructed in
Section~\ref{sec:local}.

\subsection{The density input and reduction of coefficients}

A valuation \(c\) of a function field \(B/b\), with \(b\) algebraically
closed, is a
\emph{constant reduction} if
\(\trdeg(Bc/bc)=\trdeg(B/b)\).
The composition of a constant reduction with a prime divisor of
\(Bc/bc\) is called a \emph{c.r.\ quasi-prime divisor}.
The trivial constant reduction is allowed, so ordinary prime divisors
belong to this class.
For a constant reduction, Lemma~\ref{lem:local-abhyankar} gives
\(cB=cb\): the relative value group is torsion free of rational rank
zero. After composition with a prime divisor, the relative value
group is therefore \(\Z\). The final discrete step is an innermost
convex subgroup of the composed value group. Every proper coarsening
kills this step and has all its remaining values supplied by
constants. This proves the minimality required of a quasi-prime
divisor and justifies the terminology.

For the density statement, temporarily use \(\ell\)-adic characters:
\[
 \widehat A_E=\Hom(E^\times,\Z_\ell),\qquad
 \widehat I_a=\Hom(E^\times/U_a,\Z_\ell),\qquad
 \widehat D_a=\Hom(E^\times/U_a^1,\Z_\ell).
\]
These groups have the topology of pointwise convergence. Restriction to
units and passage to residue give a map
\(\widehat\rho_v:\widehat D_v\to\widehat A_{F_0}\).
Define \(\widehat S_F^{\mathrm{in}}\) and \(\widehat J_v\) by the same
formulas as in \eqref{eq:residual-inertia-set}, with hats on the character
groups and on \(\rho_v\).

\begin{theorem}[Pop's residual density input]\label{thm:density-input}
Let \(F/\kappa\) be a function field over an algebraically closed field
of characteristic different from \(\ell\). Suppose that a valuation
\(v\) of \(F\), with \(Fv/\kappa v\) a function field, has the following
properties:
\begin{enumerate}
 \item \(vF\) has no nonzero \(\ell\)-divisible convex subgroup;
 \item there is a subfield \(\kappa\subset\kappa_1\subset F\) such that
 \[
   \kappa_1v=\kappa v,\qquad
   \trdeg(F/\kappa_1)=\trdeg(Fv/\kappa_1v)=1.
 \]
\end{enumerate}
Then \(\widehat J_v\) contains the full minimized inertia group
\(\widehat I_u\) of every c.r.\ quasi-prime divisor \(u\) of
\(Fv/\kappa v\).
\end{theorem}

This is the consequence of \cite[Theorem~1.2 and the proof of
Theorem~4.2, p.~351, item~2]{Pop2015} with the notation introduced on
p.~348. The auxiliary field \(\kappa_1\) is not required to be
algebraically closed. The minimized interpretation in residue
characteristic \(\ell\) is explained in
\cite[Remark~4.1 and the appendix, Section~A]{Pop2015}.

We verify the hypotheses for our terminal valuation \(v\).
Lemma~\ref{lem:local-flag-uniqueness} shows that \(vF\) has no nonzero
\(\ell\)-divisible convex subgroup. It remains to construct the
auxiliary field \(\kappa_1\).

Choose elements \(x_1,\ldots,x_{d-1}\in F^\times\) whose values are
rationally independent modulo \(v\kappa\), and set
\(\kappa_1=\kappa(x_1,\ldots,x_{d-1})\).
Distinct monomials in the \(x_i\) have different values modulo
\(v\kappa\). Thus every nonzero polynomial in these elements has a unique
term of least value, proving algebraic independence over \(\kappa\).
If a quotient of two such polynomials has value zero, their least terms
have the same monomial. Its residue is the residue of the quotient of
their coefficients, and consequently belongs to \(\kappa v\).
This proves \(\kappa_1v=\kappa v\); the two transcendence degrees required
by Theorem~\ref{thm:density-input} are now one.

\begin{lemma}\label{lem:residual-density-mod-ell}
For a quasi-prime \((d-1)\)-divisor \(v\) of \(F/\kappa\), the set \(J_v\)
contains the inertia line \(I_u\subset A_{F_0}\) of every c.r.\
quasi-prime divisor \(u\) of \(F_0/\kappa_0\).
\end{lemma}

\begin{proof}
Reduction of character values gives continuous maps
\[
 r_E:\widehat A_E\longrightarrow A_E.
\]
They carry \(\widehat I_a\) into \(I_a\), carry \(\widehat D_a\) into
\(D_a\), and commute with restriction to units. In particular,
\[
 r_F(\widehat S_F^{\mathrm{in}})\subset S_F^{\mathrm{in}},
 \qquad
 \rho_v r_F=r_{F_0}\widehat\rho_v
       \quad\hbox{on }\widehat D_v.
\]
The first inclusion follows from continuity and the definition as a
closure of a union.

Fix a c.r.\ quasi-prime divisor \(u\) of \(F_0/\kappa_0\). Characters with
either coefficient ring annihilate the divisible group \(u\kappa_0\).
Since \(uF_0/u\kappa_0\cong\Z\), reduction
\(\widehat I_u\to I_u\) is surjective. Given \(\chi\in I_u\), choose a
lift \(\widehat\chi\in\widehat I_u\).
Theorem~\ref{thm:density-input} supplies
\(\widehat\eta\in\widehat S_F^{\mathrm{in}}\cap\widehat D_v\) with
\(\widehat\rho_v(\widehat\eta)=\widehat\chi\).
Then \(r_F(\widehat\eta)\) belongs to \(S_F^{\mathrm{in}}\cap D_v\) and
has residual image \(\chi\). Hence \(\chi\in J_v\).
\end{proof}

The argument needs only the displayed inclusion between the two closed
unions. It does not assert that reduction commutes with their
intersection with decomposition groups, and it makes no use of a lift
of \(\Theta\).

\subsection{Valuative characters and the relation among point orders}

The density input gives a lower bound for \(J_v\). To control its other
elements, we use the fact that taking this closed union introduces no
nonvaluative characters.

\begin{lemma}\label{lem:residual-valuative}
Every element of \(S_F^{\mathrm{in}}\) is valuative. Every element of
\(J_v\) is valuative as a character of \(F_0\).
\end{lemma}

\begin{proof}
Let \(\mathcal V(F)\) be the space of valuation rings of \(F\), with its
compact patch topology. In
\(A_F\times\mathcal V(F)\), the condition that a character \(\chi\)
annihilates the unit group of a valuation ring \(O\) is closed.
Indeed, for each \(x\in F^\times\), failure of this condition is
witnessed by the open conditions
\[
 x,x^{-1}\in O,\qquad \chi(x)\ne0.
\]
The projection of the closed incidence set to \(A_F\) is compact,
hence closed. This projection is precisely the set of valuative
characters, which contains every \(I_w\) and therefore contains
\(S_F^{\mathrm{in}}\).

Now let \(\chi\in I_u\cap D_v\) for some valuation \(u\).
If \(u\leq v\), then \(\chi\in I_v\) and its residual image is zero.
If \(v\leq u\), then restriction to \(U_v\) shows that
\(\rho_v(\chi)\) annihilates the units of the induced valuation \(u/v\)
on \(F_0\).
In the remaining case let \(a\) be the finest common coarsening of
\(u\) and \(v\). Approximation for the independent induced valuations
on \(Fa\) gives
\[
 U_a=U_u U_v^1.
\]
For example, a prescribed nonzero residue can be written as an
induced \(u\)-unit times an induced principal \(v\)-unit by choosing
an element close to \(1\) at the first valuation and close to that
residue at the second. Lifting gives the displayed identity; the
remaining factor in \(U_a^1\) is already a unit for both refinements.
Thus \(\chi\) annihilates \(U_a\), and
\(\chi\in I_a\subset I_v\). Its residual image is again zero.
\end{proof}

We next record how the points of a complete curve appear in its
character space, in the finite-coefficient form of Topaz's
calculation in \cite[appendix, Lemma~5, pp.~354--355]{Pop2015}.
This description will also be used for curve subfields later.

\begin{lemma}[Point orders]\label{lem:curve-orders}
Let \(B/b\) be a function field of transcendence degree one over an
algebraically closed field, and let \(C/b\) be its smooth projective
curve. For \(x\in C(b)\), let
\(\delta_x=\ord_x\bmod\ell\in A_B\).
The lines \(\Lambda\delta_x\) are distinct and nonzero.
The family \((\delta_x)_{x\in C(b)}\) tends to zero outside finite
subsets, and the continuous summation map
\begin{equation}\label{eq:curve-summation}
 \prod_{x\in C(b)}\Lambda\longrightarrow A_B,\qquad
 (a_x)_x\longmapsto\sum_x a_x\delta_x
\end{equation}
has kernel consisting exactly of the constant families.
If \(\charac(b)\ne\ell\), the common kernel in \(V_B\) of the
\(\delta_x\) has dimension \(2g(C)\).
\end{lemma}

\begin{proof}
Approximation at distinct point valuations proves that each
\(\delta_x\) is nonzero and that their lines are distinct.
A function on \(C\) has only finitely many zeros and poles. Hence every
fixed element of \(B^\times\) is annihilated by all but finitely many
\(\delta_x\), which proves the asserted convergence and defines
\eqref{eq:curve-summation}, with continuous evaluation at each function.

A family \((a_x)_x\) defines a homomorphism from the divisor group of
\(C\) to \(\Lambda\), sending the point divisor \(x\) to \(a_x\).
It lies in the kernel of \eqref{eq:curve-summation} precisely when
this homomorphism vanishes on principal divisors, or equivalently
factors through \(\Pic(C)\).
The group \(\Pic^0(C)(b)\) is \(\ell\)-divisible: multiplication by
\(\ell\) on the Jacobian is a surjective isogeny, and \(b\) is
algebraically closed. This remains true in characteristic \(\ell\).
Every homomorphism \(\Pic(C)\to\Lambda\) therefore factors through
the degree map \(\Pic(C)\to\Z\). Since all point divisors have degree
one, its coefficient family is constant. Conversely, constant
families annihilate principal divisors.

Finally, if all point orders of \(f\in B^\times\) are divisible by
\(\ell\), write \(\operatorname{div}(f)=\ell D\) and associate to
\([f]\) the class of \(D\) in \(\Pic(C)[\ell]\). This gives an
isomorphism from the common kernel to \(\Pic(C)[\ell]\):
surjectivity follows from the definition of torsion in \(\Pic(C)\);
injectivity follows because \(b^\times\) is \(\ell\)-divisible.
When \(\charac(b)\ne\ell\), the \(\ell\)-torsion of the Jacobian is
isomorphic to \((\Z/\ell)^{2g(C)}\).
\end{proof}

\subsection{An intrinsic test for finite-field constants}

We now adapt the curve-like relation criterion of
\cite[Theorem~4.2 and appendix, Main Theorem~I, p.~353]{Pop2015}
to the family of lines in \(J_v\).
Let \(\mathcal L_v\) be the family of one-dimensional subspaces of
\(A_{F_0}\) contained in \(J_v\). Choose a nonzero generator
\(e_H\in H\) for each \(H\in\mathcal L_v\). Consider the following
conditions:
\begin{enumerate}
 \item the family \((e_H)_{H\in\mathcal L_v}\) tends to zero outside
 finite subsets;
 \item the resulting summation map
 \[
   \prod_{H\in\mathcal L_v}\Lambda\longrightarrow A_{F_0},
   \qquad (a_H)_H\longmapsto\sum_H a_H e_H
 \]
 has a one-dimensional kernel generated by a family whose every
 coordinate is nonzero.
\end{enumerate}
These conditions are independent of the choices of generators.
Changing generators rescales the coordinates of the product, and,
because \(\Lambda\) is finite, preserves the convergence in the first
condition. They are also preserved by topological linear
isomorphisms of the residual character spaces.

\begin{proposition}\label{prop:point-recognition}
For a quasi-prime \((d-1)\)-divisor \(v\) of \(F/\kappa\), the family
\(\mathcal L_v\) satisfies the preceding two conditions if and only if
\(\kappa v\) is algebraic over a finite field.
When this holds, \(\mathcal L_v\) consists exactly of the point inertia
lines of the smooth projective curve of \(Fv/\kappa v\).
Consequently the matching in Corollary~\ref{cor:flag-matching}
recognizes these terminal valuations and matches their full point
families.
\end{proposition}

\begin{proof}
Suppose first that \(\kappa_0\) is algebraic over a finite field.
Every valuation on \(\kappa_0\) is trivial. A nontrivial valuation
of the one-variable field \(F_0\) trivial on \(\kappa_0\) is a point
valuation on its smooth projective curve: properness gives a closed
center, and the discrete valuation ring at that smooth point is
dominated by the valuation ring. Writing each function as a power
of a uniformizer times a local unit identifies the two valuations.
Thus
Lemma~\ref{lem:residual-valuative} shows that every line in
\(\mathcal L_v\) is a point line. The converse inclusion follows
from Lemma~\ref{lem:residual-density-mod-ell}, using the trivial
constant reduction. Lemma~\ref{lem:curve-orders} now proves the two
conditions.

Suppose instead that \(\kappa_0\) is not algebraic over a finite field.
There is a nontrivial valuation \(b\) on \(\kappa_0\): in
characteristic zero extend a nontrivial valuation of the prime
field; in positive characteristic use a transcendental element,
extend its variable valuation to a rational transcendence basis,
and prolong to \(\kappa_0\).
Choose a transcendental \(t\in F_0\) with \(F_0/\kappa_0(t)\) finite.
The Gauss extension of \(b\) to \(\kappa_0(t)\), followed by a
prolongation to \(F_0\), gives a constant reduction \(c\).
Indeed the rational residue field has transcendence degree one,
and finite prolongation gives a finite residue extension.
The value group of \(c\) equals its constant-value subgroup: the
Gauss extension has this property, and a finite prolongation has
finite value-group index, whereas the constant-value group is
divisible.

Choose a point valuation \(q\) on \(F_0c/\kappa_0c\), and form
\(u=q\circ c\). This is a c.r.\ quasi-prime divisor, whose value
group contains the final discrete step as an innermost convex
subgroup. In particular \(uF_0/u\kappa_0\cong\Z\), so \(I_u\) is a
nonzero line, and it belongs to \(\mathcal L_v\).
It is different from every true point line. To see this, let \(p\)
be a point valuation of \(F_0/\kappa_0\). The valuation \(u\) is
nontrivial on constants, so cannot coarsen \(p\). Conversely, a
proper coarsening of \(u\) kills its innermost discrete subgroup
and retains only values supplied by constants; it cannot equal
the nontrivial valuation \(p\) that is trivial on constants.
Thus \(p\) and \(u\) are incomparable. Since \(p\) has rank one,
their finest common coarsening is trivial. Approximation gives
\(I_p\cap I_u=0\), proving the assertion even with
\(\Lambda\)-coefficients.

All point lines already belong to \(\mathcal L_v\), and their
generators have the nonzero relation of
Lemma~\ref{lem:curve-orders}, after rescaling to the chosen
generators. If the first condition holds for the full family,
extend this relation by zero on the additional lines. It is a
nonzero kernel vector with a zero coordinate, contradicting the
second condition. If the first condition fails, the family also
fails the stated test. This proves the equivalence.
\end{proof}
